% 原书第22章思考题22.1--22.12独立概念解，源书图22.30--22.32需精确重画。
\section*{第22章\quad 光的干涉：思考题解析}
\paragraph{思考题 22.1}
红与绿滤片使两缝透射场的中心光频率不同，
不同光频率的交叉项长期平均为零，
普通光探测器不能观察\emph{稳定}的
光强双缝条纹：
\[
\langle I\rangle=I_{\rm red}+I_{\rm green}
+2\Re\langle E_{\rm red}E_{\rm green}^*\rangle
\simeq I_{\rm red}+I_{\rm green}.
\]
故
\(\boxed{\text{一般看不到稳定的红绿互相干涉条纹}}\)。
若滤片谱带有明显重叠，某一\emph{共同频率成分}
仍可在相同频段产生低可见度条纹，
并非滤片颜色不同就绝对没有任何交叉项。

\paragraph{思考题 22.2}
激光输出通常具有较窄谱宽、
较高空间相干度，经双缝较容易出现高可见度条纹。
钠灯具有特征窄谱线，可提供一定时间相干长度，
但若源发光区域较宽，两缝不同位置接收到的
光场空间相关性可能较低而使条纹模糊。
通过前置窄缝或针孔选取足够小的源角宽，
钠灯\emph{同样可以}获得双缝干涉条纹。
\[
\boxed{\text{不是“钠灯绝不能干涉、激光一定干涉”；
必须比较缝距与光源相干角宽和光程差。}}
\]

\paragraph{思考题 22.3}
令两个同频相干点源 \(S_1,S_2\)
到观察点P的距离分别为 \(r_1,r_2\)，
相位差
\(\delta=k(r_1-r_2)+\delta_0\)。
同一极大或极小序号要求
\[
\boxed{r_1-r_2=\text{常数}}.
\]
几何上在平面中到两定点\emph{距离差恒定}
的轨迹为双曲线（极端零差时为垂直平分线），
因此条纹呈一族双曲线；
若两源频率不同或水速不均匀则不能简单套用。

\paragraph{思考题 22.4}
Fresnel双棱镜的两半略微偏折原始狭缝源S
发出的同一波前，使两组折射光的
反向延长线分别汇聚于
\emph{两个靠近的虚像} \(S_1,S_2\)。
它们来自同一初始光源的相干分割，
相对相位可保持稳定。
两组折射束在双棱镜后方的\emph{共同照明交叠区}
发生干涉，而不是虚像本身发射新的光能。
图22.31的几何需用Snell定律和棱镜小角
偏折近似 \(\delta\approx(n-1)A\) 定量描绘。

\paragraph{思考题 22.5}
两束同偏振、同频且相位差 \(\delta\) 恒定的光，
\[
\boxed{I=I_1+I_2+2\sqrt{I_1I_2}\cos\delta}.
\]
更一般地有复相干度 \(\gamma_{12}\)，
交叉项
\(2\sqrt{I_1I_2}\Re(\gamma_{12}e^{i\delta})\)。
若二者不相干、交叉项时间平均为零，
\[
\boxed{\langle I\rangle=I_1+I_2}.
\]
若两束\emph{正交偏振}，即使来自同一激光器，
直接不经过分析器的普通强度测量也不出现
这一交叉项；相干和偏振是两个条件。

\paragraph{思考题 22.6}
小角近似Young条纹间距
\[
\boxed{\Delta x=\frac{\lambda D}{d}}.
\]
(1) 缝距 \(d\) 增大，条纹间距减小，
相同屏宽内条纹更多。
(2) 横移原光源改变到两个缝的初始
光程差，从而移动整族条纹的\emph{中心位置}；
若光源仅改变入射角且光程几何保持均匀，
条纹间距近似不变。
(3) 发光源面积/角宽变大，
不同源点各自产生互相错开的条纹图案，
探测器的非相干叠加使条纹对比度降低，
严重时条纹消失。
在源已充分空间相干的范围内，
“光源变亮”本身不必降低可见度。

\paragraph{思考题 22.7}
同一光束分成两路并延迟 \(\tau\)
产生的干涉衡量\emph{时间相干性}：
相干时间 \(\tau_c\) 约为场的一阶相关
\(G^{(1)}(\tau)\) 显著衰减的时间尺度，
相干长度
\[
\boxed{\ell_c\sim c\tau_c}
\]
（在介质中需用适当群速）。
同一时刻不同横向位置的场互相关
\(G^{(1)}(\boldsymbol r_1,\boldsymbol r_2)\)
刻画\emph{空间相干性}；
相干宽度是场能保持显著相干的横向分离尺度。
\[
\boxed{\text{时间相干关乎光程延迟；
空间相干关乎波前上的两点间距。}}
\]

\paragraph{思考题 22.8}
肥皂泡薄膜两界面的反射光相互干涉。
膜厚 \(h\) 随位置变化，
光程差约 \(2nh\cos r\)，
加上某界面反射的\(\pi\)相位跃变，
使不同波长分别增强或抵消，
出现彩色等厚干涉纹。
空气—肥皂液—空气的反射通常只有
\emph{一个}界面发生由低到高折射率的
\(\pi\)相位跃变；
当 \(h\to0\) 时，
两反射光相差约\(\pi\)，
可见反射趋向\emph{相消而呈黑膜}。
\[
\boxed{\text{黑膜是极薄膜的反射相消，
通常是接近破裂的征兆。}}
\]
实际还需考虑表面张力、白光波段及吸收。

\paragraph{思考题 22.9}
空气劈尖小角近似
\(h(x)\simeq\alpha x+h_0\)，
等厚条纹满足
\[
\boxed{2n h(x)=m\lambda+\text{固定反射相位修正}}.
\]
(1) 整体增加厚度
\(h_0\mapsto h_0+\Delta h\)，
同序号 \(m\) 条纹满足
\(\Delta x=-\Delta h/\alpha\)；
即\emph{向原先薄端/接触端移动}。
(2) 增大劈尖角 \(\alpha\)，
\[
\boxed{\Delta x_{\rm fringe}\simeq
\frac{\lambda}{2n\alpha}}
\]
减小，条纹更密。
若整体平移同时改变楔角，需分开分析两效应。

\paragraph{思考题 22.10}
两面虽接近平行，但厚玻璃
两条反射光的光程差
\(\Delta\sim2nh\cos r\) 可能很大。
若照明只有有限时间相干长度 \(\ell_c\)，
而 \(\Delta\gg\ell_c\)，
两束光场的时间相干度接近零，
不再出现可见的明暗干涉条纹。
此外若厚板的反射光束空间分离超过探测器
重叠范围，干涉也被削弱。
\[
\boxed{\text{楔角小不等于相干光程差一定小。}}
\]

\paragraph{思考题 22.11}
超长基线射电干涉测量（VLBI）
以相距向量 \(\boldsymbol B\) 的
同步射电望远镜观察同一天体，
远场波前到达两站的几何时间差约为
\[
\boxed{\Delta t_{\rm geo}=
\frac{\boldsymbol B\cdot\hat{\boldsymbol s}}c}
\]
（符号取决于基线方向）。
与高精度原子时钟、地球自转模型及天体方位联合拟合，
可用群时延/相位时延测得基线矢量随时间的变化。
同一测站随地球自转使\(\boldsymbol B\)
在天球方向上的投影改变；
长期监测基线长度、方向改变可测
地壳板块的厘米到毫米级位移趋势。
必须校正电离层、对流层、天体结构和时钟偏差；
不是仅用一张静态干涉条纹照片就得到板块速率。

\paragraph{思考题 22.12}
薄云母片带来上方光路附加光程
\((n-1)t\)，双缝几何光程差在小角时
约 \(dy/D\)。
故
\[
\boxed{y_0=\frac{(n-1)tD}{d}}
\]
为零级中心在屏上的平移量，
\emph{向贴云母片的上缝一侧移动}。
条纹间距仍然
\[
\boxed{\Delta y=\frac{\lambda D}{d}}
\]
（在云母片不显著改变波束方向、色散和照明宽度时）。
\emph{附加常数光程差}只能改整个条纹族的位置，
不能改相邻整数级之间的间距。
